\subsection{定积分的方向、运算与检验}

前述分割求和给出了定积分。它在 $v$--$t$ 图上表示有向面积：横轴下方的部分给负位移，而路程应对速率 $|v|$ 积分。类似地，
密度不均匀细棒的质量为 $m=\int\lambda(x)\dd x$，变力做功为
$W=\int\vec F\cdot\dd\vec r$，积分的共同思想都是“局部近似简单、整体取极限求和”。

若 $F'(x)=f(x)$，则
\[
 \int_a^b f(x)\,\dd x=F(b)-F(a).
\]
积分上限、下限和积分变量缺一不可；换元时微分和边界必须同时变换。分部积分
\[
 \int_a^b u\,\dd v=[uv]_a^b-\int_a^b v\,\dd u
\]
来自乘积求导法则，适于含对数、反三角函数或两个不同类型函数乘积的积分。结果求出后可通过
求导核对原函数，并用量纲和边界极限检查物理答案。

\subsection{积分建立物理量的空间分布}

连续分布的质心坐标满足
\[
 x_C=\frac1M\int x\,\dd m.
\]
关键是先把质量元写成合适坐标的函数，如细线 $\dd m=\lambda\dd l$、薄片
$\dd m=\sigma\dd S$、立体 $\dd m=\rho\dd V$。转动惯量
$I=\int r_\perp^2\dd m$ 同样依赖转轴到质量元的垂直距离。选择对称坐标后，许多奇函数部分积分为零；
在结果中保留这种对称性，比从一开始展开全部坐标更简洁。

微分方程把“局部变化规律”还原成运动。若 $m\dot v=F(v)$，可写成
$\dd v/F(v)=\dd t/m$ 再积分；更一般的 $F(v,t)$ 只有在它能够分解为速度函数与时间函数的乘积时才能分离变量。
由 $v\dd v/\dd x=a$ 可在不显含时间的问题中把二阶运动方程降阶。
积分常数由初始条件确定，不能因只求到一个原函数而省略。遇到分段力或分段约束时，应在每一段
分别积分，并在转换时刻用位置、速度的连续性或冲量条件衔接。

积分也能给出曲线长度、面积和体积。平面极坐标的一般面积元为
$\dd S=r\dd r\dd\theta$，柱坐标体积元为 $\dd V=r\dd r\dd\theta\dd z$；这些微元中的几何因子
不是记号装饰，而是坐标网格随位置伸缩的结果。若区域从原点延伸到边界 $r=R(\theta)$，先对半径积分才得到
$\int_0^{R(\theta)}r\dd r\dd\theta=\tfrac12R^2(\theta)\dd\theta$；后者是微小扇形面积，不是一般面积元。
先画出被积区域、再确定积分次序，能够避免漏算或重复。

\subsection{近似求积与物理积分的检查}

有些积分没有初等原函数，但定积分仍然有清楚的物理意义。可以把区间划分为小段，用矩形、梯形或
中点值近似每段面积；分割越细，数值和越接近真实积分。实验数据本来就是离散的，此时
\[
 \int_{t_1}^{t_2}F(t)\,\dd t
 \simeq\sum_i\frac{F_i+F_{i+1}}2(t_{i+1}-t_i)
\]
可直接由测得的力—时间数据估计冲量。数值积分给出的是近似值，应同时说明时间步长和误差怎样随
分割改变，而不能把计算器输出的许多小数位都当成有效数字。

广义积分还要求检查端点或无穷远是否收敛。例如总质量
$M=\int_0^\infty\lambda(x)\dd x$ 只有在密度衰减足够快时才有限；势能或电势把无穷远取作零点时，
也隐含了相应积分收敛。若积分发散，往往表示理想化模型不能延伸到无限远，而不是简单地说物理量
“等于无穷”。无限长导线、无限大平面等模型通常只能可靠给出场强或两个有限点之间的势差。

含参数的积分能够表示一族物理状态。若积分区间固定并满足适当光滑条件，可在积分号下求导：
\[
 \frac{\dd}{\dd a}\int_{x_1}^{x_2}f(x,a)\,\dd x
 =\int_{x_1}^{x_2}\frac{\partial f}{\partial a}\,\dd x.
\]
这个方法可把难积分转化为较容易的积分，再对参数积分回来；积分常数则由某个易求的参数值确定。
若上下限也依赖参数，还要加上端点移动产生的边界项。

任何物理积分至少应通过三项检查：被积式乘微元的量纲是否正确，改变积分方向时符号是否合理，
以及在均匀分布、恒力或恒速极限下能否退化为熟悉的初等公式。先写清微元代表什么，再写积分限，
往往比先套积分表更能避免错误。
